\documentclass{report}
\usepackage{amsmath,amssymb}
\setlength{\parindent}{0mm}
\setlength{\parskip}{1em}
\begin{document}
\begin{center}
\rule{6in}{1pt} \
{\large Andrew Canning \\
{\bf Hybrid OpenMP/MPI Conjugate Gradient Iterative Eigensolver for First-principles Plane Wave Electronic Structure Codes on Many-core Architectures }}

LBNL MS-50F \\ One Cyclotron Road \\ Berkeley \\ CA 94720 \\ USA
\\
{\tt acanning@lbl.gov}\end{center}

Density Functional Theory (DFT) based First-principles materials science
codes using plane waves
(PW Fourier basis) have become the largest user, by method, of computer
cycles at scientific computer centers around
the world. At NERSC (National Energy Research Scientific Computing
Center) an estimated 17% of the total
cycles are used by DFT-PW codes such as VASP, Quantum Espresso, Abinit,
PEtot, PARATEC etc.
These codes commonly use conjugate gradient based iterative eigensolvers
to solve the density functional theory based
approximation to the many-body Schrodinger Equation (usually the Kohn-Sham form).
In this approach 3D FFTs are used to move between real and Fourier space
to construct the matrix-vector product
such that the different parts of the Hamiltonian matrix are calculated in
the space where they are sparse.
The parallel scaling of the 3D FFTs in the conjugate gradient solver is
particularly challenging as rather than one large grid,
as is the case in other scientific applications using spectral methods,
we have many medium sized grids
(one for each electronic state) which can limit scaling. We therefore
developed a specialized hybrid OpenMP/MPI
version of the 3D FFT to scale efficiently on modern many-core platforms.

Overall the conjugate gradient solver uses a two level parallelization
scheme where the high level parallelization
divides the eigenvectors (electronic states) among groups of nodes and
then within a group each eigenvector is
divided among the nodes. OpenMP and threading is then used to parallelize
the solver on the node while MPI is
used for the communications between the nodes. Details of how OpenMP and
threading is used to parallelize the 3D FFT
and other parts of the solver will be given in the talk.
We will present results for the complete code as well as separately for
the 3D FFT on the many-core architecture
Intel MIC Xeon NERSC computers Edison (Cray XC30 with 12 core Intel Ivy
Bridge Xeon) and Cori phase one (Cray with 16 core
Intel Haswell Xeon).
By sending fewer larger messages the hybrid OpenMP/MPI
version of the 3D FFT significantly outperforms the pure MPI version on
large core count many-core architectures
allowing the solver to scale efficiently to 10,000s of cores. This work
was done in collaboration with L-W Wang, J. Shalf,
N.J. Wright (LBNL), M. Gajbe (NCSA) and S. Anderson (Cray Inc.).
This research used resources of the National Energy Research Scientific
Computing Center, which is supported by
the Office of Science of the U.S. Department of Energy under Contract No.
DE-AC02-05CH11231.


\end{document}
